% =========================================================================
%  Figura: propiedades básicas. (a) Conductor cargado: campo normal a la
%  superficie, de módulo sigma/eps0, y camino interior entre A y B
%  (volumen equipotencial). (b) Cilindro de Gauss en la superficie.
%  Contorno en polares: r(t) = 1.45 + 0.30 cos 2t + 0.18 sen 3t; la normal
%  exterior se calcula con la derivada r'(t) en cada punto.
% =========================================================================
\begin{figure}[t]
\centering
\begin{tikzpicture}[font=\small]
  % ---------------------------------------------------------------- (a)
  \begin{scope}
    \fill[top color=azulvelo!80!white, bottom color=metal, draw=azulprof, line width=0.6pt]
      plot[domain=0:360, samples=120, smooth cycle, variable=\t]
      ({(1.45+0.30*cos(2*\t)+0.18*sin(3*\t))*cos(\t)},{(1.45+0.30*cos(2*\t)+0.18*sin(3*\t))*sin(\t)});
    \foreach \t in {0,20,...,340}{
      \pgfmathsetmacro{\r}{1.45+0.30*cos(2*\t)+0.18*sin(3*\t)}
      \pgfmathsetmacro{\rp}{-0.60*sin(2*\t)+0.54*cos(3*\t)}
      \pgfmathsetmacro{\tx}{\rp*cos(\t)-\r*sin(\t)}
      \pgfmathsetmacro{\ty}{\rp*sin(\t)+\r*cos(\t)}
      \pgfmathsetmacro{\nn}{sqrt(\tx*\tx+\ty*\ty)}
      \pgfmathsetmacro{\nx}{\ty/\nn}\pgfmathsetmacro{\ny}{-\tx/\nn}
      \draw[-{Stealth[length=4pt]}, azulacero, line width=0.6pt]
        ({\r*cos(\t)},{\r*sin(\t)}) -- ++({0.5*\nx},{0.5*\ny});
      \node[signopos, font=\tiny\bfseries] at ({\r*cos(\t)-0.13*\nx},{\r*sin(\t)-0.13*\ny}) {$+$};}
    % camino interior entre A y B
    \draw[cargapos!80, line width=0.7pt, densely dashed, -{Stealth[length=4pt]}]
      (-0.85,0.05) .. controls (-0.3,1.0) and (0.4,-0.5) .. (0.75,0.45);
    \fill[azulprof] (-0.85,0.05) circle (1.3pt) node[left, font=\footnotesize] {$A$};
    \fill[azulprof] (0.75,0.45) circle (1.3pt) node[right, font=\footnotesize] {$B$};
    \node[font=\footnotesize, text=azulprof, align=center] at (0.1,-0.75) {$\vE=\vc{0}$,\ $\phi=\mathrm{cte}$};
    \node[font=\footnotesize, text=azulacero] at (2.35,1.25) {$\vE=\dfrac{\sigma}{\eo}\,\un$};
    \node[font=\footnotesize, align=center] at (0,-2.55)
      {$\displaystyle\phi_B-\phi_A=-\int_A^B\vE\cdot\dl=0$};
    \node at (0,-3.25) {(a)};
  \end{scope}
  % ---------------------------------------------------------------- (b)
  \begin{scope}[xshift=6.6cm, yshift=-0.2cm]
    \fill[top color=azulvelo!80!white, bottom color=metal] (-2.2,-1.6) rectangle (2.2,0);
    \draw[azulprof, line width=0.7pt] (-2.2,0) -- (2.2,0);
    \foreach \x in {-1.9,-1.5,...,1.95}{\node[signopos, font=\tiny\bfseries] at (\x,-0.12) {$+$};}
    % cilindro de Gauss
    \draw[gaussazul, fill=white, fill opacity=0.35] (-0.75,-0.75) rectangle (0.75,0.75);
    \draw[gaussazul] (0,0.75) ellipse (0.75 and 0.13);
    \draw[gaussazul] (0,-0.75) ellipse (0.75 and 0.13);
    \foreach \x in {-0.45,0,0.45}{\draw[-{Stealth[length=4pt]}, azulacero, line width=0.7pt] (\x,0.05) -- (\x,1.55);}
    \draw[-{Stealth[length=4pt]}, verdeazul, line width=0.8pt] (1.15,0.75) -- (1.15,1.35) node[right, font=\footnotesize] {$\un$};
    \node[font=\footnotesize, text=azulacero] at (0.0,1.85) {$\vE=E\,\un$};
    \node[font=\footnotesize] at (1.3,0.42) {$A$};
    \node[font=\footnotesize, text=azulprof] at (-1.45,-1.05) {$\vE=\vc{0}$};
    \node[font=\footnotesize, text=grispizarra] at (1.7,-1.25) {conductor};
    \node[font=\footnotesize, text=grispizarra] at (-1.65,0.55) {vacío};
    \node[font=\footnotesize, align=center] at (0,-2.35)
      {$\displaystyle\oiint\vE\cdot\dS=EA=\frac{\sigma A}{\eo}$};
    \node at (0,-3.05) {(b)};
  \end{scope}
\end{tikzpicture}
\caption{Propiedades básicas del conductor en equilibrio. (a) Conductor cargado: el interior es un volumen equipotencial sin campo (camino
de $A$ a $B$ por dentro del metal), la carga está en la superficie y el campo exterior es normal a ella, de módulo $\sigma/\eo$. (b) Cálculo del
campo superficial con un cilindro de Gauss: solo la base exterior contribuye al flujo.}
\label{fig:propiedades}
\end{figure}
